

\chapter{Formal Design with PlusCal and TLA+}
\section{TLA+ and PlusCal in a Nutshell}
\begin{itemize}
    \item State machines, variables, actions.
    \item PlusCal syntax: labels, \texttt{either/or}, \texttt{assert}, invariants.
    \item Safety (bad states never occur) vs.\ liveness (good states eventually occur) --- the two obligations a TM model must state separately.
    \item Why a model is a design tool, not just a proof tool: the act of writing the model forces the empty corner cases into view.
\end{itemize}

\section{Specifying and Checking Simple TM Algorithms}
\begin{itemize}
    \item Single‑global‑lock TM in PlusCal.
    \item NOrec (value‑based validation, timestamp‑based writes) in PlusCal.
    \item Translating PlusCal to TLA+ and running TLC.
    \item The companion repository contains checked models for every backend in this book: SGL, TL2, NOrec, TinySTM variants, MVLog, JVSTM, plus GPU models (PR-STM, JVSTM, GUST) at warp granularity.
    \item Invariants that matter: opacity/reduced order, lock exclusion, commit-clock monotonicity, ``no committed transaction can be invalidated.''
    \item TLC as an oracle that catches the ``obviously correct'' bug --- e.g.\ the read-validate\slash{}snapshot-threshold bugs discussed in Part~IV and Part~VII.
\end{itemize}

\begin{table}[htbp]
\centering
\small
\begin{tabular}{p{0.15\linewidth}p{0.21\linewidth}p{0.28\linewidth}p{0.26\linewidth}}
\hline
\textbf{Design family} & \textbf{Representative backend} & \textbf{Modelled shared state} & \textbf{Main safety check} \\
\hline
Serial lock & SGL, PersistentSGL & lock bit, heap & mutual exclusion, opacity \\
Value validation & NOrec, NOrec-BF & sequence clock, read values & stale read rejection \\
Versioned locks & TL2, TSC-TM & orecs, versions, ownership & lock/version consistency \\
Undo/redo STM & TinySTM variants & ownership, write sets & ownership before write-back \\
Multi-version & JVSTM, MVLog & versions, commit order & visible snapshot consistency \\
GPU commit-log & GUST & log, CCT, MRV tables & no missed conflict \\
Distributed 2PC & TiKV/Percolator & locks, writes, timestamps & atomic visibility \\
\hline
\end{tabular}
\caption{A small taxonomy of what the model must remember. The runtime code may contain allocators, barriers, persistence fences, or RPCs; the model keeps only the state needed to falsify the advertised TM property.}
\label{tab:tla-taxonomy}
\end{table}

\section{Worked Example: A Single-Global-Lock TM in PlusCal}
\index{PlusCal}\index{TLA+}\index{TLC}\index{invariant}\index{model checking}\index{counterexample}\index{liveness}
\label{sec:pluscal-example}

The simplest TM in this book, in the notation that Chapter~6's invariants
check.  Figure~\ref{fig:pluscal-latex} shows the same algorithm twice:
once as the PlusCal source that is model-checked (left), and once as the
pseudo-code that \textsc{TLATeX} (\texttt{tla2tex.TLA}) typesets from that
same source (right).  This is the point of PlusCal: the algorithm in a
paper and the specification checked by TLC are \emph{the same text}.

\begin{figure}[ht]
\centering
\begin{minipage}[t]{0.47\textwidth}
\centering\textbf{(a) PlusCal source}\\[2pt]
\begin{lstlisting}[basicstyle=\ttfamily\footnotesize]
--algorithm sgl
variables lock = FALSE, x = 0, y = 0;
process Tx \in 1..2
variables rs = {}, ws = {};
begin
  begin_tx:
    await lock = FALSE; lock := TRUE;
  body:
    rs := rs \cup {x, y};          (* reads *)
    ws := ws \cup {x, y};          (* writes *)
  commit:
    x := 2; y := 2;                (* write-back *)
    lock := FALSE;
end process
\end{lstlisting}
\end{minipage}\hfill
\begin{minipage}[t]{0.47\textwidth}
\centering\textbf{(b) \textsc{TLATeX} pseudo-code output}\\[2pt]
{\footnotesize
\begin{tabbing}
  \hspace{0.9cm}\=\hspace{3.1cm}\=\hspace{2.6cm}\=\hspace{2.5cm}\=\kill
  \textbf{--algorithm} \emph{sgl} \\
  \textbf{variables} \> \texttt{lock = FALSE, x = 0, y = 0}; \\
  \textbf{process} \> \texttt{Tx in 1..2} \\
  \> \textbf{variables} \> \texttt{rs = \{\}, ws = \{\}}; \\
  \textbf{begin} \\
  \> \texttt{begin\_tx:} \\
  \> \> \textbf{await} \> \texttt{lock = FALSE}; \texttt{lock := TRUE}; \\
  \> \texttt{body:} \\
  \> \> \texttt{rs := rs union \{x, y\}}; \> \> reads \\
  \> \> \texttt{ws := ws union \{x, y\}}; \> \> writes \\
  \> \texttt{commit:} \\
  \> \> \texttt{x := 2; y := 2}; \> \> write-back \\
  \> \> \texttt{lock := FALSE}; \\
  \textbf{end process}
\end{tabbing}
}
\end{minipage}
\caption{A single-global-lock TM, twice. (a) is the PlusCal source kept in the repository and model-checked with TLC; (b) is the pseudo-code that \textsc{TLATeX} produces from the \emph{same} file --- the paper figure can never drift from the verified algorithm.}
\label{fig:pluscal-latex}
\end{figure}

\begin{itemize}
    \item Every invariant of this chapter is trivially true here --- the lock serialises everything --- which is exactly why SGL is the reference design every optimistic backend is checked against.
    \item The companion repository's \texttt{SGL.tla} is this model plus the opacity invariant, and it passes TLC.
    \item Exercise~1 asks you to add read-set validation to this model; the point is that even this ``obviously correct'' design needs the check spelled out.
    \item Contrast with the optimistic designs of Part~IV: the moment the lock disappears, the invariants become non-trivial and TLC earns its keep.
\end{itemize}

\section{Bug‑Injection and Counterexample Analysis}
\begin{itemize}
    \item Deliberately omitting a read‑set validation.
    \item Reading the TLC counterexample trace to localise the bug.
    \item Mutation testing for models: remove the guard, remove the fence, remove the re-validation, and confirm TLC finds each injected bug.
    \item Counterexamples vs.\ crashes: a TLC trace before an invariant violation is a precise, replayable schedule --- exactly the artifact a debugger cannot easily produce for real runs.
\end{itemize}

\section{Worked Example: Bank Transfer Under a Single Global Lock}
\index{bank transfer}\index{single global lock}\index{PlusCal}

The running bank example is a two-account transfer.  The model is intentionally
small: two accounts, two transactions, one unit transferred from account~1 to
account~2.  The interesting part is not the arithmetic; it is that the floor
and conservation checks are written at the same point where the runtime would
return success to the caller.

\begin{lstlisting}[caption={A complete PlusCal model of the bank transfer under SGL.}]
---- MODULE BankSGL ----
EXTENDS Naturals, TLC

Accounts == 1..2
Initial == [a \in Accounts |-> 10]
Total(b) == b[1] + b[2]

(*--algorithm BankSGL
variables
  lock = FALSE,
  bal = Initial,
  committed = [t \in 1..2 |-> FALSE];

process Tx \in 1..2
variables
  from = 1,
  to = 2,
  amount = 1,
  seen = 0;
begin
BeginTx:
  await lock = FALSE;
  lock := TRUE;

Read:
  seen := bal[from];

WriteBack:
  if seen >= amount then
    bal[from] := bal[from] - amount;
    bal[to] := bal[to] + amount;
    committed[self] := TRUE;
  end if;

Check:
  assert Total(bal) = 20;
  assert bal[1] >= 0 /\ bal[2] >= 0;

Commit:
  lock := FALSE;
end process;
end algorithm; *)
====
\end{lstlisting}

\begin{itemize}
    \item The invariant is checked after write-back, not only at initialization.
    \item The lock makes the state space small enough that TLC explores all interleavings of both transactions.
    \item The model does not claim performance.  It claims that every checked execution preserves the account floors and total money.
    \item SGL is a useful reference because every optimistic backend should refine the same externally visible transfer behavior.
\end{itemize}

\section{Worked Example: NOrec-Style Validation for the Bank}
\index{NOrec}\index{validation}\index{opacity}\index{read set}

NOrec replaces ownership metadata with a global sequence clock and value-based
validation.  A reader records a snapshot and the values it observed.  If the
clock changes before commit, the read values must still match memory; otherwise
the transaction aborts.  The following model is deliberately finite and omits
allocation, privatization, and signal recovery.  It is the core validation
discipline.

\begin{lstlisting}[caption={A finite PlusCal model of NOrec-style value validation.}]
---- MODULE BankNOrecValue ----
EXTENDS Naturals, TLC

Accounts == 1..2
Initial == [a \in Accounts |-> 10]
Total(b) == b[1] + b[2]

(*--algorithm BankNOrecValue
variables
  commitLock = FALSE,
  clock = 0,
  bal = Initial;

process Tx \in 1..2
variables
  snap = 0,
  r1 = 0,
  r2 = 0,
  amount = 1,
  ok = FALSE;
begin
BeginTx:
  snap := clock;

Read1:
  r1 := bal[1];

Read2:
  r2 := bal[2];

Validate:
  if clock = snap then
    ok := TRUE;
  else
    ok := (r1 = bal[1] /\ r2 = bal[2]);
  end if;

Commit:
  if ok /\ r1 >= amount then
    await commitLock = FALSE;
    commitLock := TRUE;

    if clock = snap /\ r1 = bal[1] /\ r2 = bal[2] then
      bal[1] := bal[1] - amount;
      bal[2] := bal[2] + amount;
      clock := clock + 1;
    end if;

    assert Total(bal) = 20;
    assert bal[1] >= 0 /\ bal[2] >= 0;

    commitLock := FALSE;
  end if;
end process;
end algorithm; *)
====
\end{lstlisting}

\begin{itemize}
    \item The first validation is a fast path: if the clock did not change, the read set is still valid.
    \item The second validation is the commit-time defense.  Removing it is a standard bug-injection test.
    \item Value validation compares the values read, not only the timestamp.  This is the point that distinguishes NOrec from a naive global-clock scheme.
    \item The model is small, but it captures the critical opacity question: no transaction may commit after observing a mixed snapshot.
\end{itemize}

\begin{figure}[htbp]
\centering
\resizebox{\textwidth}{!}{%
\begin{tikzpicture}[
    node distance=1.8cm and 1.9cm,
    state/.style={draw, rounded corners, minimum width=2.8cm, minimum height=0.85cm, align=center},
    edge/.style={-{Latex[length=2mm]}, thick}
]
\node[state] (begin) {Begin\\record clock};
\node[state, right=of begin] (read) {Read\\record values};
\node[state, right=of read] (validate) {Validate\\clock or values};
\node[state, right=of validate] (commit) {Commit\\lock and recheck};
\node[state, below=of validate] (abort) {Abort\\discard log};

\draw[edge] (begin) -- (read);
\draw[edge] (read) -- (validate);
\draw[edge] (validate) -- node[above, font=\footnotesize]{valid} (commit);
\draw[edge] (validate) -- node[right, font=\footnotesize]{stale} (abort);
\draw[edge] (commit) |- node[pos=0.25, right, font=\footnotesize]{failed recheck} (abort);
\draw[edge] (abort) -| ($(begin.south)+(0,-0.4)$) -- (begin.south);
\end{tikzpicture}%
}
\caption{Validation state machine for the finite NOrec-style model. TLC explores all interleavings through these states; the useful counterexamples are usually the paths where validation succeeds too early or commit omits the final recheck.}
\label{fig:norec-validation-sm}
\end{figure}

\section{The Complete Verification Workflow}
\label{sec:verify-workflow}
What ``we checked this with TLC'' means, end to end, so the claim is
reproducible rather than rhetorical:

\begin{enumerate}
    \item \textbf{Specify} the algorithm in PlusCal (a \texttt{--algorithm} block), with the correctness condition encoded as invariants (e.g.\ opacity via \texttt{InvNoMissedConflict}).
    \item \textbf{Add temporal properties} where liveness matters (a committed-transaction progress property under weak fairness).
    \item \textbf{Configure TLC}: a \texttt{.cfg} listing the model constants, the invariant, the optional \texttt{PROPERTY}, and a state bound so the checker terminates.
    \item \textbf{Run} \texttt{tlc} and read the state-count line --- \emph{how much} of the state space was explored matters as much as the absence of errors.
    \item \textbf{Inject a bug} (Section ``Bug-Injection'') and confirm the checker finds it: a test of a model that never fails is a test of the model's triviality.
    \item \textbf{Document the limits}: bounded state spaces and abstraction gaps (no call stack, no jmp-buffer lifetime) are what the model cannot see.
\end{enumerate}

\begin{table}[ht]
\centering
\small
\begin{tabular}{llll}
\hline
\textbf{Runtime (Part IV/VII)} & \textbf{TLA+ model} & \textbf{Key invariant} & \textbf{Result} \\
\hline
SGL (Ch.~8)     & \texttt{SGL.tla}         & opacity (all invariants) & TLC pass \\
NOrec (Ch.~9)   & \texttt{NOrec.tla}       & value-validation opacity  & TLC pass \\
TL2 (Ch.~9)     & \texttt{TL2.tla}         & versioned-lock correctness & TLC pass \\
TinySTM (Ch.~9) & \texttt{TinySTM\_WBCTL.tla} & lock/ownership invariants & TLC pass \\
Romulus (Ch.~9) & \texttt{Romulus.tla}     & read-validate OCC          & TLC pass \\
GPU GUST (Ch.~17)& \texttt{GPU\_GUST.tla}  & \texttt{InvNoMissedConflict} & TLC pass \\
\hline
\end{tabular}
\caption{Runtime-to-model traceability in the companion repository: every chapter that introduces a runtime names the PlusCal/TLA+ file, the invariant checked, and the outcome. The repository's \texttt{docs/proofs} directory is the source of truth; this table is its index.}
\label{tab:traceability}
\end{table}

\begin{itemize}
    \item The workflow is the connective tissue between Part~III (formalism) and Part~IV (runtimes): a backend ``is verified'' only in the sense of steps 1--6 above.
    \item Figure~\ref{fig:pluscal-latex} is step~1's artifact: the same source that TLC checks is the figure a reader sees.
    \item Figure~\ref{fig:norec-validation-sm} is the corresponding artifact for optimistic validation: it names the states where injected bugs should be placed.
\end{itemize}

\section{Summary}
\begin{itemize}
    \item Model checking provides formal guarantees about TM algorithm correctness.
    \item Even simple TM designs can harbour subtle bugs that TLC exposes.
    \item The model is not the implementation.  It is the executable argument that the implementation's concurrency protocol has the advertised shape.
    \item The useful models are small, finite, and hostile: they should make wrong interleavings easy for TLC to find.
\end{itemize}

\section{Exercises}
\begin{enumerate}
    \item Extend the single‑global‑lock PlusCal model with an explicit read‑set validation step. Use TLC to check that opacity holds.
    \item The NOrec model can fail if it uses only timestamp validation without value‑based validation. Insert this bug and show the TLC output that reveals it.
    \item Specify a TM that uses eager locking and use TLC to find a deadlock scenario.
    \item \emph{[implementation]} Add a tiny C++ harness that runs the bank transfer under the repository's SGL and NOrec backends with two worker threads. Check the same two runtime assertions used in the models: total money is conserved and no account falls below its floor.
    \item \emph{[verification]} Mutate the protocol of Figure~\ref{fig:norec-validation-sm} by removing the commit-time recheck from \texttt{BankNOrecValue}. State the invariant that fails and summarize the shortest TLC counterexample trace.
    \item \emph{[theory]} Explain why lock exclusion alone is insufficient for opacity in an optimistic STM. Give a two-transaction history with a stale read, then identify which validation edge in Figure~\ref{fig:norec-validation-sm} must reject it.
\end{enumerate}


%  PART IV: IMPLEMENTING TRANSACTIONAL MEMORY
